% ---------------------------------------------------------------------------
% Nerd Nite Fort Collins house style for beamer.
%
% Palette and typeface lifted from the real poster source,
% ~/Documents/NerdNiteFortCollins/git/NerdNiteFortCollins/POSTERS/S2025E08.tex
% Do not invent brand colours here; if the poster changes, change this.
% ---------------------------------------------------------------------------
\usepackage{fontspec}
\usepackage{tikz,graphicx,xcolor}
\usetikzlibrary{calc}

\definecolor{NNblue}{RGB}{0, 42, 55}
\definecolor{NNred}{RGB}{237, 28, 37}
\definecolor{NNLblue}{RGB}{19, 218, 249}
\definecolor{NNpink}{RGB}{248, 164, 167}
\definecolor{NNorange}{RGB}{255, 97, 27}
\definecolor{NNpaper}{RGB}{245, 243, 238}

\setsansfont[Ligatures=TeX]{Courier New}
\renewcommand{\familydefault}{\sfdefault}

\usetheme{default}
\setbeamercolor{background canvas}{bg=NNblue}
\setbeamercolor{normal text}{fg=white}
\setbeamercolor{structure}{fg=NNLblue}
\setbeamertemplate{navigation symbols}{}
\setbeamertemplate{footline}{}
\setbeamertemplate{headline}{}
\beamertemplatenavigationsymbolsempty
\setlength{\parindent}{0pt}

% Slides must never hyphenate. "wire-less" across two lines on slide 3 is how
% this was found; a broken word on a projector reads as a typo from row ten.
\hyphenpenalty=10000
\exhyphenpenalty=10000
\tolerance=2000
\emergencystretch=3em

% --- furniture -------------------------------------------------------------
% Logo bug + show mark, bottom right. On every frame but the close.
\newcommand{\bug}{%
  \begin{tikzpicture}[remember picture,overlay]
    \node[anchor=south east,inner sep=0]
      at ($(current page.south east)+(-6mm,5mm)$)
      {\includegraphics[height=7mm]{assets/NNFoCoLogo.pdf}};
    \node[anchor=south east,inner sep=0,text=NNLblue,font=\tiny]
      at ($(current page.south east)+(-16mm,8mm)$)
      {Nerd Nite Fort Collins};
  \end{tikzpicture}}

% Cyan rule under a heading, the poster's banner idea at slide scale.
\newcommand{\nnrule}{{\color{NNLblue}\rule{\linewidth}{1.2pt}}}

% Slide heading.
\newcommand{\nnhead}[1]{%
  {\color{NNLblue}\large\bfseries #1}\\[1.5mm]\nnrule\\[4mm]}

% --- image slots -----------------------------------------------------------
% \slot{name}{width}{height}{what photograph belongs here}
%
% SIZING, settled 13 Oct. The slide is 16.0 x 9.0 cm and the text block is
% ~14.5 cm, so a two-column photo box is capped by its COLUMN, not by the
% page height -- every landscape image here has 2-4 cm of unused vertical.
% Each box is therefore set to its image's own aspect ratio at the widest the
% column allows, which is why the numbers look arbitrary:
%   1.777 (16:9)  -> 6.8 x 3.83      1.333 (4:3) -> 6.8 x 5.10
%   1.641 (mars)  -> 6.8 x 4.14      1.359 (collage, wider column) -> 8.1 x 5.96
% p03-trajectories is the one exception: it is portrait (0.751) and therefore
% HEIGHT-limited, maxed at 5.6 x 7.45 with ~6 mm of page left under it.
% nntext + nnphoto deliberately sum under 1.0; at 0.97 beamer's column
% separation pushed the pair ~2 mm past the text block on every frame.
%
% Drops in img/<name>.jpg or img/<name>.png the moment the file exists.
% Until then it draws a labelled placeholder box of exactly the same size, so
% the layout you are looking at now is the layout you get with the real photo.
% \slot{name}{width}{height}{what photograph belongs here}
%
% SIZING, settled 13 Oct. The slide is 16.0 x 9.0 cm and the text block is
% ~14.5 cm, so a two-column photo box is capped by its COLUMN, not by the
% page height -- every landscape image here has 2-4 cm of unused vertical.
% Each box is therefore set to its image's own aspect ratio at the widest the
% column allows, which is why the numbers look arbitrary:
%   1.777 (16:9)  -> 6.8 x 3.83      1.333 (4:3) -> 6.8 x 5.10
%   1.641 (mars)  -> 6.8 x 4.14      1.359 (collage, wider column) -> 8.1 x 5.96
% p03-trajectories is the one exception: it is portrait (0.751) and therefore
% HEIGHT-limited, maxed at 5.6 x 7.45 with ~6 mm of page left under it.
% nntext + nnphoto deliberately sum under 1.0; at 0.97 beamer's column
% separation pushed the pair ~2 mm past the text block on every frame.
%
% WIDTH AND HEIGHT ARE BARE NUMBERS IN CM -- {6.2}{4.4}, never {6.2cm}.
% The first version took dimensions with units and then did arithmetic on
% them ("4.4cm-1.25"), which tikz accepted and placed the labels outside the
% box. The picture sets x=1cm,y=1cm so plain numbers are centimetres.
%
% Drops in img/<name>.jpg|png|jpeg the moment the file exists. Until then it
% draws a labelled box of exactly the same size, so the layout on screen now
% is the layout you get with the real photograph in it.
\newcommand{\slotreal}[3]{%
  \begin{tikzpicture}[x=1cm,y=1cm]
    \useasboundingbox (0,0) rectangle (#2,#3);
    \fill[NNblue] (0,0) rectangle (#2,#3);  % letterbox blends into the ground
    \begin{scope}
      \clip (0,0) rectangle (#2,#3);
      \node[inner sep=0] at (#2/2,#3/2)
        {\includegraphics[width=#2 cm,height=#3 cm,keepaspectratio]{#1}};
    \end{scope}
    \draw[NNLblue,line width=0.8pt] (0,0) rectangle (#2,#3);
  \end{tikzpicture}}

\newcommand{\slotph}[4]{%
  \begin{tikzpicture}[x=1cm,y=1cm]
    \useasboundingbox (0,0) rectangle (#2,#3);
    \fill[NNblue!65!black] (0,0) rectangle (#2,#3);
    \draw[NNLblue,dashed,line width=0.7pt] (0,0) rectangle (#2,#3);
    \draw[NNorange,line width=1.6pt] (0,#3-0.55) -- (0,#3) -- (0.55,#3);
    \draw[NNorange,line width=1.6pt] (#2-0.55,0) -- (#2,0) -- (#2,0.55);
    \node[anchor=north west,text=NNorange,font=\tiny,inner sep=0]
      at (0.28,#3-0.3) {PHOTO};
    \node[anchor=north east,text=NNLblue,font=\tiny,inner sep=0]
      at (#2-0.15,#3-0.3) {img/#1};
    \node[anchor=north west,text=white,font=\tiny,align=left,
          text width=#2 cm-0.6cm,inner sep=0] at (0.28,#3-0.75) {#4};
  \end{tikzpicture}}

\newcommand{\slot}[4]{%
  \IfFileExists{img/#1.jpg}{\slotreal{img/#1.jpg}{#2}{#3}}{%
   \IfFileExists{img/#1.png}{\slotreal{img/#1.png}{#2}{#3}}{%
    \IfFileExists{img/#1.jpeg}{\slotreal{img/#1.jpeg}{#2}{#3}}{%
     \slotph{#1}{#2}{#3}{#4}}}}}

% --- slide skeletons -------------------------------------------------------
% Two-column frames use beamer's columns, not minipages. The first version
% used minipages at 0.50 + 0.46 of \textwidth plus 8mm margins either side,
% which overflowed and silently dropped the photo below the text.
\newenvironment{nnframe}{%
  \begin{frame}[plain]\bug\vspace*{3mm}}{\end{frame}}

\newcommand{\nntext}{0.49\textwidth}
\newcommand{\nnphoto}{0.47\textwidth}
